\section{Completion obstructions and vector-parameter geometry}
\label{obs-section}

The positive normalized arithmetic jump energy is an exact control, but its
conversion into the Weil form requires a negative contact and a signed pole
contribution.  We first identify an obstruction to a scalar ground-state
completion.  We then construct a vector-parameter connection compatible with
the benchmark Berry curvature and show why that compatibility does not fix the
missing contact.  Neither argument excludes positive bulk theories with more
general source maps or internal degrees of freedom.

\subsection{The complete form and its scalar ground transform}

Fix $I=(-L/2,L/2)$ and a finite prime set $\mathcal S$ containing every prime
with $\log p<L$.  For zero extensions $F=E_Lf$ and $G=E_Lg$, set
\begin{gather}
 n_\gamma(r)=\frac{e^{-|r|/2}}{1-e^{-2|r|}},\qquad
 a_{p,m}=(\log p)p^{-m/2},\qquad d_{p,m}=m\log p,\nonumber\\
 \nu_{\mathcal S}(dr)=n_\gamma(r)\,dr+
 \sum_{p\in\mathcal S,m\ge1}a_{p,m}
       (\delta_{d_{p,m}}+\delta_{-d_{p,m}})(dr),\label{obs-jump-measure}\\
 j(f,g)=\frac12\int\nu_{\mathcal S}(dr)\int_\R
 \overline{F(x+r)-F(x)}\,[G(x+r)-G(x)]\,dx.\label{obs-jump-form}
\end{gather}
The gamma density is of order $1/(2|r|)$ at zero, and the prime measure has
finite mass.  Thus $j$ is the nonnegative gamma form plus a bounded
perturbation, with the same logarithmic Fourier form domain.  Write $J$ for
its associated operator.  The complete target has the exact form
\begin{equation}
 W_L=J-\kappa_{\mathcal S}I+P_L,
 \qquad \kappa_{\mathcal S}=-w_0+2\sum_{p\in\mathcal S,m\ge1}a_{p,m}>0,
 \qquad w_0=\psi(1/4)-\log\pi,
 \label{obs-full-operator}
\end{equation}
where
\begin{equation}
 p(x,y)=2\cosh((x-y)/2),\qquad
 P_L=2|c\rangle\langle c|-2|s\rangle\langle s|,
 \quad c(x)=\cosh(x/2),\quad s(x)=\sinh(x/2).
 \label{obs-pole-kernel}
\end{equation}
Although $p(x,y)$ is positive, $P_L$ is an indefinite rank-two operator.
Inactive prime jumps contribute only a diagonal term to $j$, canceled by
$\kappa_{\mathcal S}$; adding such primes does not change $W_L$.

Let $\mu_{\mathcal S}$ be the internal jump measure on $I^2$ obtained from
\eqref{obs-jump-measure}, and let
$k_{\rm esc}(x)=\int_{x+r\notin I}\nu_{\mathcal S}(dr)$.
Then $j$ is the internal difference form plus
$\int_I k_{\rm esc}\overline f g$.
Take a real $h\in C^1(\overline I)$ bounded above and bounded away from zero.
All identities below hold first for $f=hu$, $g=hv$ with
$u,v\in C_c^\infty(I)$ and with zero exterior values for $h$ in $Jh$.
Expansion and symmetrization give
\begin{equation}
 j(hu,hv)=\frac12\int_{I^2}h(x)h(y)
 (\overline{u(x)-u(y)})(v(x)-v(y))\,d\mu_{\mathcal S}
 +\int_I h(Jh)\overline u v.
 \label{obs-ground-identity}
\end{equation}
Consequently the potential $-(Jh)/h$ admits a positive incidence-square
representation.  This is the useful scalar ground-state mechanism: a
negative potential can arise from a positive energy after a change of
ground state.  For example, $h=1$ removes the positive exterior escape
rate.  Multiplying $h$ by a constant, however, does not change $(Jh)/h$.
The normalization of a vacuum ray cannot select the prescribed contact.
The incidence operator can be closed to define a positive form; identifying
that closure with a prescribed operator sum is a further domain question.
Our obstructions use only the displayed compactly supported core.

The pole term has the corresponding identity
\begin{equation}
\begin{aligned}
 \ip{f}{P_Lg}={}&\int_I r_P(x)\overline f g\\
 &-\frac12\int_{I^2}p(x,y)
       (\overline{f(x)-f(y)})(g(x)-g(y))\,dx\,dy,\\
 r_P(x)={}&8\sinh(L/4)\cosh(x/2).
\end{aligned}
 \label{obs-pole-difference}
\end{equation}
Combining the two expansions yields the full polarized ground identity
\begin{equation}
 \begin{split}
 q_L(hu,hv)={}&\frac12\int_{I^2}h(x)h(y)
 (\overline{u(x)-u(y)})(v(x)-v(y))
 [d\mu_{\mathcal S}-p(x,y)\,dx\,dy]\\
 &+\int_I h(W_Lh)\overline u v.
 \end{split}
 \label{obs-signed-ground-identity}
\end{equation}
Even a positive solution of $W_Lh=0$ would leave a signed edge measure.
The poles change off-diagonal couplings and therefore cannot be treated
merely as a local killing potential.

\subsection{A sign obstruction before the first prime}

\begin{proposition}[Scalar source obstruction]
\label{obs-scalar-proposition}
Let $\rho>1$ be the root of $\rho^3-\rho-1=0$ and put $r_*=\log\rho$.
For every $L>r_*$ there are nonzero nonnegative smooth inputs with disjoint
supports for which $q_L(u,v)>0$.  Therefore, with a pointwise scalar source
identification, $q_L$ cannot be a positive-ground-state transform of a real
scalar diffusion or jump form with nonnegative edge weights, even after
adding an arbitrary local potential.
\end{proposition}

\begin{proof}
Away from the finitely many active prime delays, the off-diagonal kernel is
\begin{equation}
 R(r)=2\cosh(r/2)-n_\gamma(r)
 =e^{-r/2}\frac{e^{3r}-e^r-1}{e^{2r}-1},\qquad r>0.
 \label{obs-sign-kernel}
\end{equation}
The polynomial $x^3-x-1$ is strictly increasing for $x>1$, so $R$ is
positive precisely when $r>r_*$.  Choose $d\in(r_*,L)$ away from all
active prime delays and two sufficiently narrow patches separated by $d$.
Their difference set avoids those delays and stays above $r_*$.  For
nonnegative $u,v$ supported in the respective patches,
$q_L(u,v)=\int u(x)R(|x-y|)v(y)\,dx\,dy>0$.
The contact and all local potentials vanish on these disjoint inputs.
In contrast, a nonnegative scalar jump measure has a nonpositive cross
term on disjoint nonnegative inputs; a local diffusion contributes zero.
Positive multiplication by $h$ preserves both disjointness and sign.
\end{proof}

The restriction concerns a Markov-type scalar boundary energy, not
positivity of an arbitrary Hilbert-space quadratic form.  It also persists
under scalar reversible bulk elimination.  At a finite regulator a positive
real symmetric diffusion/jump matrix has nonpositive off-diagonal entries.
Its invertible interior block has entrywise nonnegative inverse, so the
Schur complement again has nonpositive off-diagonal entries.  Limits
preserving the bilinear form on these test inputs retain the same sign.

\begin{proposition}[Scalar phase obstruction]
\label{obs-phase-proposition}
For $L>2r_*$, a scalar phase change cannot turn all the off-diagonal
couplings into real nonpositive Markov couplings.  This obstruction already
occurs for some $L<\log2$.
\end{proposition}

\begin{proof}
Choose three narrow patches with all pairwise separations above $r_*$ and
away from prime delays.  Their continuous couplings are positive.  A phase
change would have to assign phase $-1$ to each of the three edges.
The product of phases around a triangle is $1$ for a scalar gauge, whereas
the required product is $(-1)^3=-1$.  The argument can be formulated almost
everywhere on the patches and does not require pointwise regularity of the
phase.  Finally, evaluating $x^3-x-1$ at $\sqrt2$ shows
$\rho<\sqrt2$, hence $2r_*<\log2$.
\end{proof}

Matrix states, genuine magnetic holonomy, and nonlocal or derivative source
maps are outside these exclusions.  The triangle identifies why scalar
rephasing is insufficient; it does not exclude an energy containing
interference squares with positive cross terms.

\subsection{The compulsory diagonal and the contact problem}

Indeed the pole kernel has a positive interference-square construction:
\begin{equation}
 e_P^+(f,g)=\frac12\int_{I^2}p(x,y)
 (\overline{f(x)+f(y)})(g(x)+g(y))\,dx\,dy
 =\ip{f}{P_Lg}+\int_I r_P\overline f g.
 \label{obs-positive-pole-square}
\end{equation}
The equality follows simply by expanding the square.  Its consequence is
\begin{equation}
 q_L[f]=j[f]+e_P^+[f]
       -\int_I(\kappa_{\mathcal S}+r_P(x))|f(x)|^2\,dx.
 \label{obs-compulsory-diagonal}
\end{equation}
Thus a positive square reproduces the signed pole operator only together
with an additional diagonal energy.  The completion must derive the
subtraction of this companion as well as the original contact.  Adding two
pole amplitudes does not by itself perform that task.

More precisely, put $D=\kappa_{\mathcal S}I-P_L$.  This bounded operator has
infinite rank: it equals $\kappa_{\mathcal S}I$ on
$\{c,s\}^{\perp}$.  A finite-rank correction of the prepared state map
cannot remove the entire difference $J-W_L$.  This does not exclude a
finite number of bulk field species, each carrying infinitely many modes.

The danger of a circular Schur construction is already explicit in a
valid one-prime support window.  If $\log2<L\le\log3$, then $D>0$:
$c\perp s$,
$\norm{c}^2=L/2+\sinh(L/2)$, and
$\norm{s}^2=-L/2+\sinh(L/2)$, so its smallest possible exceptional
eigenvalue is $\kappa_{\mathcal S}-2\norm{c}^2>0$.
Here $2\norm{c}^2\le\log3+2/\sqrt3<3$, whereas
$\kappa_{\mathcal S}\ge-w_0>3$.
The elementary square completion
\begin{equation}
 \mathcal A[f,z]=j[f]-2\re\ip{D^{1/2}f}{z}+\norm{z}^2
 =q_L[f]+\norm{z-D^{1/2}f}^2
 \label{obs-circular-schur}
\end{equation}
does have the desired minimized value.  But $\mathcal A\ge0$ holds if and
only if $q_L\ge0$.  Its positivity is exactly the unknown assertion.
Adding $\ip{f}{Df}$ makes the action manifestly positive and changes its
minimum to $j[f]$.  A useful coherent constraint must therefore supply an
independently derived infinite-dimensional contraction or projection
identity, including the pole sectors, rather than repackage the target's
unknown positivity.

\subsection{A vector-parameter lift and its normalization freedom}

Vector-multiplet parameters offer a different compatibility equation from
ordinary chiral $tt^*$ geometry.  Four-supercharge invariance constrains
their Berry connection and projected moment-map operator by a Bogomolny
equation \cite{SonnerTong}.  This is a necessary effective equation, not a
converse construction of a microscopic theory.

For the finite-prime benchmark define, on $\re z>0$,
\begin{equation}
 \Lambda_{\mathcal S}(z)=\pi^{-z/2}\Gamma(z/2)
       \prod_{p\in\mathcal S}(1-p^{-z})^{-1},\qquad
 \ell(z)=\log\Lambda_{\mathcal S}(z).
 \label{obs-local-factor}
\end{equation}
The logarithm exists and is real on the positive axis.  The specified
ground family has $\norm{\Psi_\sigma}^2=e^{\ell(\sigma)}$ and normalized
extension $\Omega_{\sigma,\eta}=e^{i\eta X/2}\Omega_\sigma$.
In the convention $A=i\ip{\Omega}{d\Omega}$,
\begin{equation}
 A_\sigma=0,\qquad A_\eta=-\tfrac12\ell'(\sigma),\qquad
 F_{\sigma\eta}=-\tfrac12\ell''(\sigma),\qquad
 \ell''(\sigma)=\operatorname{Var}_\sigma(X)>0.
 \label{obs-real-berry}
\end{equation}
Introduce an auxiliary third coordinate $v$, orient parameter space as
$(\sigma,\eta,v)$, and define
\begin{equation}
 A_\sigma=A_v=0,\qquad
 A_\eta=-\tfrac12\re\ell'(\sigma+iv),\qquad
 \Phi=-\tfrac12\im\ell'(\sigma+iv).
 \label{obs-bogomolny-lift}
\end{equation}
The Cauchy--Riemann equations give
$\nabla\times A=\nabla\Phi$: its $\sigma$ and $v$ components are,
respectively, $-\im\ell''/2$ and $-\re\ell''/2$ on both sides.
At $v=0$ this restricts to \eqref{obs-real-berry}, with $\Phi=0$.
This explicit half-space solution works for any suitable holomorphic
$\ell$ real on the positive axis.  It neither selects arithmetic factors
nor makes the auxiliary extension the actual Berry data of a quantum model.

For completeness,
\begin{equation}
 \ell'(z)=-\frac{\log\pi}{2}+\frac12\psi(z/2)
       -\sum_{p\in\mathcal S}\frac{\log p}{p^z-1}.
 \label{obs-log-derivative}
\end{equation}
Its continued local-factor poles have logarithmic-derivative residue $-1$,
with multiplicities at coincidences.  A simple boundary pole produces a
two-dimensional dipole-type harmonic $\Phi$, invariant along $\eta$,
rather than a regular three-dimensional point monopole.  Global defects
and boundary conditions would require separate physical specifications.

The contact ambiguity remains exact.  Replacing $\ell(z)$ by
$\ell(z)+cz$, $c\in\R$, corresponds to rescaling the original source
$\Psi_\zeta$ by $e^{c\zeta/2}$, where $\zeta=\sigma+i\eta$ labels that
chosen complex family.  Its normalized state changes only by
$e^{ic\eta/2}$.  Consequently $A_\eta$ shifts by $-c/2$, a flat pure gauge
on the noncompact $\eta$-line, while the curvature and $\Phi$ stay fixed.
Nevertheless
\begin{equation}
 \partial_\sigma\log|\Lambda_{\mathcal S}(\sigma+i\tau)|^2
 \longmapsto
 \partial_\sigma\log|\Lambda_{\mathcal S}(\sigma+i\tau)|^2+2c.
 \label{obs-contact-gauge}
\end{equation}
Thus curvature cannot determine this scalar contact.  Changing the origin
of $X$ changes its mean but not its variance; a fixed physical source must
supply the missing convention.

At $\sigma=1/2$ the pole-free arithmetic multiplier is
$2\re\ell'(1/2+i\tau)=-4A_\eta(1/2,\tau)$.
This equality identifies functions, not a positive norm.  The choice
$v=\tau$ is an additional proposed correspondence, and the signed pole
form is still absent.  Scalar completion factors do not restore that form
by a pointwise critical-line derivative.  A successful enlargement must
derive a physical source and the complete pairing, including its contact
and residue or boundary contributions.  These requirements leave matrix,
fermionic, magnetic, and relative-boundary mechanisms open while specifying
the precise failure of the scalar and curvature-only repairs tested here.
